\section{Algorithm Design}


\subsection{Population Initialization}

The initial population combines:

\begin{enumerate}
\item \textbf{Seed campaigns}: Human-designed campaigns providing domain knowledge.
\item \textbf{Random variants}: Randomly generated campaigns for diversity.
\item \textbf{Historical winners}: Top performers from previous campaigns.
\end{enumerate}

\begin{algorithm}
\caption{Population Initialization}
\begin{algorithmic}[1]
\REQUIRE Seeds $S$, population size $N$
\STATE $P \leftarrow S$
\STATE $P \leftarrow P \cup \text{loadHistorical}(|S|)$
\WHILE{$|P| < N$}
    \STATE $G \leftarrow \text{randomGenome}()$
    \IF{$\text{isValid}(G)$}
        \STATE $P \leftarrow P \cup \{G\}$
    \ENDIF
\ENDWHILE
\RETURN $P$
\end{algorithmic}
\end{algorithm}

\subsection{Selection}

We employ tournament selection with elitism:

\begin{algorithm}
\caption{Tournament Selection}
\begin{algorithmic}[1]
\REQUIRE Population $P$, tournament size $k$, elite count $e$
\STATE $\text{elite} \leftarrow \text{top}_e(P, f)$
\STATE $\text{selected} \leftarrow \text{elite}$
\WHILE{$|\text{selected}| < |P|$}
    \STATE $T \leftarrow \text{sample}(P, k)$
    \STATE $\text{winner} \leftarrow \arg\max_{G \in T} f(G)$
    \STATE $\text{selected} \leftarrow \text{selected} \cup \{\text{winner}\}$
\ENDWHILE
\RETURN selected
\end{algorithmic}
\end{algorithm}

Elitism preserves the top $e$ individuals unchanged, preventing regression.

\subsection{Crossover Operators}

We define domain-specific crossover operators:

\subsubsection{Creative Crossover}

For text genes (headline, body), we use semantic-aware crossover:

\begin{lstlisting}[language=Python,caption=Semantic text crossover]
def crossover_text(parent1, parent2):
    # Parse into semantic units
    units1 = parse_semantic_units(parent1)
    units2 = parse_semantic_units(parent2)

    # Combine units preserving grammar
    child_units = []
    for i, (u1, u2) in enumerate(zip(units1, units2)):
        if random() < 0.5:
            child_units.append(u1)
        else:
            child_units.append(u2)

    return reconstruct(child_units)
\end{lstlisting}

Semantic units include noun phrases, verb phrases, and adjective modifiers. Crossover combines units while preserving grammatical validity.

\subsubsection{Image Crossover}

For image genes, we cannot directly combine pixels. Instead, we crossover image metadata:

\begin{lstlisting}[language=Python,caption=Image metadata crossover]
def crossover_image(parent1, parent2):
    # Crossover style attributes
    child_style = {
        "brightness": choice([parent1, parent2]).brightness,
        "contrast": choice([parent1, parent2]).contrast,
        "saturation": choice([parent1, parent2]).saturation,
    }

    # Select base image
    base = choice([parent1, parent2]).base_image

    return apply_style(base, child_style)
\end{lstlisting}

\subsubsection{Targeting Crossover}

Audience targeting uses set-based crossover:

\begin{lstlisting}[language=Python,caption=Audience crossover]
def crossover_targeting(parent1, parent2):
    child = {}

    # Numeric ranges: interpolate
    child["age_min"] = (parent1.age_min + parent2.age_min) // 2
    child["age_max"] = (parent1.age_max + parent2.age_max) // 2

    # Set attributes: union with probability
    child["interests"] = set()
    for interest in parent1.interests | parent2.interests:
        if random() < 0.7:  # Bias toward inclusion
            child["interests"].add(interest)

    return child
\end{lstlisting}

\subsection{Mutation Operators}

Mutation introduces variation through domain-specific operators:

\subsubsection{Text Mutation}

\begin{lstlisting}[language=Python,caption=Text mutation operators]
def mutate_text(text, rate):
    if random() < rate:
        op = choice(["synonym", "rephrase", "shorten", "expand"])
        if op == "synonym":
            return replace_with_synonym(text)
        elif op == "rephrase":
            return rephrase_sentence(text)
        elif op == "shorten":
            return remove_filler_words(text)
        else:
            return add_power_words(text)
    return text
\end{lstlisting}

Synonym replacement uses WordNet \cite{wordnet1995}. Rephrasing uses a sequence-to-sequence model trained on paraphrase corpora.

\subsubsection{Targeting Mutation}

\begin{lstlisting}[language=Python,caption=Targeting mutation]
def mutate_targeting(targeting, rate):
    if random() < rate:
        op = choice(["narrow", "broaden", "shift"])
        if op == "narrow":
            targeting.interests = random_subset(targeting.interests, 0.8)
        elif op == "broaden":
            related = get_related_interests(targeting.interests)
            targeting.interests |= random_subset(related, 0.2)
        else:
            targeting.age_min += randint(-5, 5)
            targeting.age_max += randint(-5, 5)
    return targeting
\end{lstlisting}

\subsection{Fitness Evaluation}

Fitness evaluation requires deploying campaigns and measuring results. This presents challenges:

\begin{enumerate}
\item \textbf{Cost}: Each evaluation costs real advertising spend.
\item \textbf{Latency}: Conversions may occur hours or days after impression.
\item \textbf{Noise}: Small sample sizes yield high-variance estimates.
\end{enumerate}

We address these through:

\textbf{Surrogate modeling.} A neural network predicts fitness from genome features, reducing required evaluations:

\begin{equation}
\hat{f}(G) = \text{NN}(\phi(G))
\end{equation}

where $\phi(G)$ extracts features (text embeddings, audience size, historical performance of similar campaigns).

\textbf{Bayesian updating.} We maintain posterior distributions over conversion rates:

\begin{equation}
p(\theta | D) \propto p(D | \theta) p(\theta)
\end{equation}

where $\theta$ is the true conversion rate and $D$ is observed data. We use Beta-Binomial conjugate priors.

\textbf{Budget allocation.} We allocate budget to individuals proportional to their information value:

\begin{equation}
\text{budget}_i = \text{budget}_{\text{total}} \cdot \frac{\sigma_i^2}{\sum_j \sigma_j^2}
\end{equation}

where $\sigma_i^2$ is the variance of individual $i$'s fitness estimate.

\subsection{Termination Criteria}

The algorithm terminates when:

\begin{enumerate}
\item Budget exhausted: $\sum_G \text{spend}(G) \geq \text{budget}$.
\item Convergence: Fitness improvement below threshold for $k$ generations.
\item Time limit: Campaign end date reached.
\end{enumerate}

